\documentclass[12pt]{scrartcl}

% LaTeX Template für Abgaben an der Universität Stuttgart
% Autor: Sandro Speth
% Forked Version: JSendra
%-----------------------------------------------------------
% Modul fuer verwendete Pakete.
% Neue Pakete einfach einfuegen mit dem \usepackage Befehl:
% \usepackage[options]{packagename}

% Encoding and Fonts
\usepackage[utf8]{inputenc}
\usepackage[T1]{fontenc}
\usepackage[ngerman]{babel}
\usepackage{lmodern}

% Graphics and Colors
\usepackage{graphicx}
\usepackage[pdftex,hyperref,dvipsnames,table]{xcolor}
\usepackage{tikz}
\usetikzlibrary{automata,positioning,arrows.meta}

% Math packages
\usepackage{amsmath,amssymb,amstext,amsthm}

% Lists, URLs, etc.
\usepackage[inline]{enumitem} % Ermöglicht ändern der enum Item Zahlen
\usepackage{hyperref}
\usepackage{url}
\usepackage{booktabs}
\usepackage{float}
\usepackage{pifont}
\usepackage{tcolorbox}

% Plotting
\usepackage{pgfplots}

% Code Listings
\usepackage{listings}
\usepackage{minted}

% UML Class Diagrams
% One Day I may fix the UML Diagrams
%\usepackage{styles/tikz-uml}
%\usetikzlibrary{positioning}

% Algorithms
\usepackage[lined,algonl,boxed]{algorithm2e}
% alternative zu algorithm2e:
%\usepackage[]{algorithm} %counter mit chapter
%\usepackage{algpseudocode}

% Geometry
\usepackage[a4paper,lmargin={2cm},rmargin={2cm},tmargin={3.5cm},bmargin = {2.5cm},headheight = {4cm}]{geometry}

% Headers
\usepackage[headsepline]{scrlayer-scrpage} 
\pagestyle{scrheadings} 

\pgfplotsset{compat=1.18}


% LaTeX Template für Abgaben an der Universität Stuttgart
% Autor: Sandro Speth
% Bei Fragen: Sandro.Speth@iste.uni-stuttgart.de
% Contributor: JSendra - extended and modified for current Abgaben schema
%-----------------------------------------------------------
% Modul beinhaltet Befehl fuer Aufgabennummerierung,
% sowie die Header Informationen.

% Überschreibt enumerate Befehl, sodass 1. Ebene Items mit
\renewcommand{\theenumi}{(\alph{enumi})}
% (a), (b), etc. nummeriert werden.
\renewcommand{\labelenumi}{\text{\theenumi}}

% Counter für das Blatt und die Aufgabennummer.
% Ersetze die Nummer des Übungsblattes und die Nummer der Aufgabe
% den Anforderungen entsprechend.
% Gesetz werden die counter in der hauptdatei, damit siese hier nicht jedes mal verändert werden muss
% Beachte:
% \setcounter{countername}{number}: Legt den Wert des Counters fest
% \stepcounter{countername}: Erhöht den Wert des Counters um 1.
\newcounter{sheetnr}
\newcounter{exnum}

\newcounter{teilaufgabe}[exnum]


% Befehl für die Aufgabentitel
\newcommand{\exercise}[1]{
    \section*{
        Aufgabe \theexnum\stepcounter{exnum}: #1
        }
} % Befehl für Aufgabentitel

\newcommand{\teil}[1]{
    \section*{
        Aufgabe \theexnum.\theteilaufgabe\stepcounter{teilaufgabe}: #1
    }
}


% Formatierung der Kopfzeile
% \ohead: Setzt rechten Teil der Kopfzeile mit
% Namen und Matrikelnummern aller Bearbeiter
\ohead{Josse Sendra (3542251)\\
Dennis Hehn (3593682)\\
Matthias Wagner (3542811)}
% \chead{} kann mittleren Kopfzeilen Teil sezten
% \ihead: Setzt linken Teil der Kopfzeile mit
% Modulnamen, Semester und Übungsblattnummer
\ihead{NumStatStoch\\
Sommersemester 26\\
Blatt \thesheetnr}




% Java Style lstlisting

\lstdefinestyle{javaStyle}{
            language=Java,
            backgroundcolor=\color{gray!5},
            basicstyle=\ttfamily\small,
            keywordstyle=\color{blue!80!black}\bfseries,
            stringstyle=\color{teal!70!black},
            commentstyle=\color{gray!80}\itshape,
            numbers=left,
            numberstyle=\tiny\color{gray},
            stepnumber=1,
            frame=single,
            framerule=0.5pt,
            breaklines=true,
            showstringspaces=false,
            tabsize=2,
            morekeywords={var},
            xleftmargin=0pt,
            framexleftmargin=0pt
        }

\lstnewenvironment{javaSnippet}{
  \lstset{style=javaStyle}
}{}


% New list Style, i. ii. iii. ...
\newlist{enumRoman}{enumerate}{1}
\setlist[enumRoman]{label=\roman*.}

% New list Style, a) b) ...
\newlist{enumAlph}{enumerate}{1}
\setlist[enumAlph]{label=\alph*)}

% Java Console
    % Colors

\definecolor{consolebg}{RGB}{245,245,245}   % light gray background, like printed output
\definecolor{consoletext}{RGB}{0,0,0}      % black text

\renewcommand{\ttdefault}{lmtt} % Latin Modern Typewriter, looks nicer than default

    % Define style
    
\lstdefinestyle{javaOutputStyle}{
    basicstyle=\ttfamily\small\color{consoletext}, % monospaced text
    backgroundcolor=\color{consolebg},            % light gray background
    breaklines=true,
    showstringspaces=false,
    frame=none,                                   % NO horizontal lines
    xleftmargin=0pt,
    xrightmargin=0pt,
    aboveskip=1em,
    belowskip=1em,
}

\lstnewenvironment{javaOutput}{
    \lstset{style=javaOutputStyle}
}{}

\SetKwInOut{Input}{Input}
\SetKwInOut{Output}{Output}


\setcounter{sheetnr}{4}
\setcounter{exnum}{4}
\setcounter{teilaufgabe}{1}


\begin{document}

% Aufgabe 4.1 
\teil{Komplexe Wiederholung, DFT}

\centering{$n:=4$ und $\omega_n := e^{2\pi/n}$}
\begin{enumAlph}

    \item $k=0,1,2,3$,  $m_k:=(\omega_n^0,\omega_n^k,\omega_n^{2k},\omega_n^{3k})^T \in \mathbb{C}
    ^4$\\
    
    Mit $n=4$ sind die Werte von $\omega_4$ wie folgt:

    \[\omega_n = \omega_4 = e^{2\pi i/4}=e^{i\pi /2}=i\]
    \[\omega_4^{0}=1, \omega_4^{1}=i,\omega_4^{2}=-1,\omega_4^{3}=-i\]

    Und die zugehörige Vektoren:\\
    
    \vspace{0.1cm}
\begin{center}%
    $\mathit{m_0} = \begin{pmatrix} 1 \\ 1 \\ 1 \\ 1 \end{pmatrix}$% 
    $\mathit{m_1} = \begin{pmatrix} 1 \\ i \\ -1 \\ -i \end{pmatrix}$%
    $\mathit{m_2} = \begin{pmatrix} 1 \\ -1 \\ 1 \\ -1 \end{pmatrix}$%
    $\mathit{m_3} = \begin{pmatrix} 1 \\ -i \\ -1 \\ i \end{pmatrix}$%
\end{center}

   \item Matrizen $M$ und $\overline{M}=\overline{(\omega_4^{jk})}_{0\leq j,k\leq 3}=\omega_4^{-jk}$\\

   \vspace{0.2cm}

     \[   
     M=\begin{pmatrix}
     1 & 1 & 1 & 1      \\
     1 & i & -1 & -i    \\
     1 & -1 & 1 & -1    \\
     1 & 1 & -1 & i     \\
     \end{pmatrix}
    ,
    \overline{M}=\begin{pmatrix}
     1 & 1 & 1 & 1      \\
     1 & -i & -1 & i    \\
     1 & -1 & 1 & -1    \\
     1 & i & -1 & -i     \\
     \end{pmatrix}
    \]    
    \item Matrizenprodukt $M\overline{M}$\\
    
    Das Matrizenprodukt ist als 
    $(M\overline{M})_{j,l} = \sum_{k=0}^{3} \omega_4^{jk}\cdot \overline{\omega_4^{lk}}
    = \sum_{k=0}^{3} \omega_4^{jk}\cdot \omega_4^{-lk}
    = \sum_{k=0}^{3} \omega_4^{(j-l)k}$
    definiert. Hier folgt eine Fallunterscheidung.

    \begin{itemize}
        \item[] Fall 1: $j = l$
        In diesem Fall werden alle Terme $\omega_4^0 = 1$ und deswegen ist die Summe gleich 4.

        \item[] Fall 2: $j \neq l$
        Die Summe ist eine Geometrische Reihe und da $\omega_4^4=1$ ist, folgt:

        \[\sum_{k=0}^{3}(\omega_4^{(j-l)})^k = \frac{(\omega_4^{(j-l)})^4-1}{\omega_4^{j-l}-1}=\frac{1-1}{\omega_4^{j-l}-1}=0\]
        
    \end{itemize}

    Daraus folgt:
        
    \[
    M\overline{M} = 4 \cdot I_4=
    \begin{pmatrix}
    4 & 0 & 0 & 0      \\
    0 & 4 & 0 & 0    \\
    0 & 0 & 4 & 0    \\
    0 & 0 & 0 & 4     \\
    \end{pmatrix}
    \]
        
    \newpage
    \item 
    Der $(j,l)$-Eintrag des Produkts $M\overline{M}$ lässt sich wie folgt schreiben:

    \[(M\overline{M})_{j,l}=\sum_{k=0}^{3} \omega_4^{jk}\cdot \overline{\omega_4^{lk}}
    = \sum_{k=0}^{3} \omega_4^{jk}\cdot \omega_4^{-lk}= \langle m_j,m_l \rangle\]\\

    Da wir in c) schon gezeigt haben, dass $M\overline{M} = 4I_4$. Es ergeben sich wieder 2 Fälle:

    \begin{itemize}
        \item[] Für $j \neq l$ ist $\langle m_j,m_l \rangle = 0$. Daraus folgt, dass die Vektoren $m_0,m_1,m_2$ und $m_3$ daher paarweise orthogonal sind.

        \item[] Für $j=l$ ist $\langle m_j,m_l \rangle = \lVert m_j \lVert ^2 = 4$. Somit ist $\lVert m_j \lVert = 2$ für alle $j=0,1,2,3$.
        
    \end{itemize}

    
     
\end{enumAlph}


\newpage



% Aufgabe 4.2
\teil{Polynommultiplikation mittels Fouriertransformation}

\begin{enumAlph}
    \item 
    \[N = N_u + N_v - 1\] 
    N wird gewählt, weil, wenn $u(x)$ den Grad $N_u - 1$ und $v(x)$ den Grad $N_v - 1$ hat, ihr Produkt den Grad $N_u + N_v - 2$ hat und daher genau $N_u + N_v - 1$ Koeffizienten benötigt.

    Wird $N$ kleiner als $N_u + N_v - 1$ gewählt, berechnet die DFT eine zyklische statt einer regulären Faltung, was zu Aliasing führt, sodass das Polynomprodukt fehlerhaft wird.

    Wird $N$ größer gewählt, funktioniert der Algorithmus zwar weiterhin korrekt, führt jedoch unnötige zusätzliche Berechnungen durch und ist daher weniger effizient.

    \item Die ,,geeigneten Stellen'' \\
    \[\omega_k = e^{2\pi i k / N}, k=0,...,N-1\]
    sind komplexe Zahlen, die gleichmäßig auf dem Einheitskreis verteilt sind. Die DFT berechnet das Polynom auf diesen Punkten.

    For jedes $k$,\\
    \[u(\omega_k) = \sum_{j=0}^{N-1} u_j\omega_k^j\]
    \[v(\omega_k) = \sum_{j=0}^{N-1} v_j\omega_k^j\]

    Diese Punkte werden gewählt, weil sie eine effiziente Auswertung und Interpolation mittels der FFT ermöglichen und weil ein Polynom vom Grad höchstens $N-1$ durch seine Werte an $N$ verschiedenen Punkten eindeutig bestimmt ist. 

    \item Für das Produktpolynom 
    \[ \omega(x) = u(x)v(x)\]
    gilt an jedem Bewertungspunkt
    \[w(\omega_k) = u(\omega_k)v(\omega_k) \]
    Die komponentenweise Multiplikation der bewerteten Werte ergibt somit die Werte des Produktpolynoms.
    Da ein Polynom vom Grad höchstent $N-1$ durch diese $N$ Werte eindeutig bestimmt ist, rekonstruiert die Anwendung der inversen DFT die Koeffizienten
    \[(w_0,...,w_{N-1})\]
    die genau die Koeffizienten des Produkpolynoms sind.

    \item Die direkte Berechnung über
    \[w_i = \sum_{j+k=i} u_j v_k\]
    entspricht einer diskreten Faltung und benötigt quadratisch viele Operationen:
    $\mathcal{O}(N^2)$
    Das beschriebene Verfahren verwendet zwei FFTs, eine komponentenweise Multiplikation und eine inverse FFT. Da jede FFT Aufwand
    $\mathcal{O}(N\log N)$
    hat, ergibt sich insgesamt
    $\mathcal{O}(N\log N)$
    Damit ist das FFT-Verfahren deutlich effizienter als die direkte Berechnung.
\end{enumAlph}

\end{document}